% ============================================================================
%  Chapter 3 --- State of the Art
% ============================================================================
\chapter{State of the Art}
\label{ch:stateofart}

\section{Introduction}

This chapter surveys the current state of AI-driven telecom operations, anomaly detection techniques, and OSS--BSS convergence approaches. We position our work relative to existing academic and industrial solutions, identifying the gaps that our AI Operations Agent addresses.

\section{AI in Telecommunications Operations}

\subsection{The AIOps Paradigm}

AIOps (Artificial Intelligence for IT Operations) has emerged as the dominant framework for applying machine learning to operational data. Originally coined by Gartner, AIOps platforms ingest telemetry data from multiple sources, apply ML models for anomaly detection, event correlation, and root cause analysis, and produce actionable recommendations.

In the telecom domain, AIOps extends to:
\begin{itemize}
  \item \textbf{Network performance prediction:} Using historical KPI data to forecast SLA breaches before they occur.
  \item \textbf{Anomaly detection:} Identifying abnormal network behaviour from cell-level metrics without explicit labels.
  \item \textbf{Root cause analysis:} Correlating symptoms across multiple network elements to locate the origin of degradation.
  \item \textbf{Self-healing:} Automatic remediation actions triggered by AI decisions (ADN L4/L5).
\end{itemize}

\subsection{Existing Approaches}

Table~\ref{tab:related-work} summarises relevant approaches in the literature.

\begin{table}[H]
\centering
\caption{Comparison of related AI-driven telecom approaches.}
\label{tab:related-work}
\small
\begin{tabularx}{\textwidth}{@{}lXcccc@{}}
\toprule
\textbf{Approach} & \textbf{Focus} & \textbf{Real Data} & \textbf{O+B} & \textbf{Cloud} & \textbf{Agent} \\
\midrule
SELFNET~\cite{selfnet2017} & Self-healing 5G networks & Simulated & No & No & No \\
ETSI ZSM~\cite{etsi_zsm} & Zero-touch management architecture & Spec only & Partial & Yes & Partial \\
Nokia AVA~\cite{nokia_ava} & Telecom analytics platform & Yes & No & Yes & No \\
Huawei SmartCare~\cite{smartcare} & CEM with demarcation & Yes & Partial & Yes & No \\
Ericsson EOM~\cite{ericsson_eom} & Expert-system operations & Yes & No & Yes & No \\
Academic IF-based~\cite{if_telecom} & Anomaly detection on KPIs & Synthetic & No & No & No \\
\textbf{Our Project} & \textbf{CEM--CVM AI Agent} & \textbf{Yes (TT)} & \textbf{Yes} & \textbf{Yes} & \textbf{Yes} \\
\bottomrule
\end{tabularx}
\end{table}

The key observation is that existing solutions either focus on the OSS side (network anomaly detection) or the BSS side (revenue analytics), but rarely bridge both. When they do (SmartCare, ETSI ZSM), the intelligence layer is either manual or limited to rule-based correlation. Our project is unique in combining real data, O+B convergence, cloud-native deployment, and an agentic AI architecture.

\section{Anomaly Detection in Telecom Networks}

\subsection{Supervised vs. Unsupervised Approaches}

Anomaly detection in telecommunications can be approached through supervised or unsupervised methods:

\begin{itemize}
  \item \textbf{Supervised:} Requires labelled data (normal vs. anomalous). Common classifiers include Random Forest, XGBoost, and neural networks. The challenge is obtaining labelled data in production environments --- operators rarely label every anomaly event comprehensively.
  \item \textbf{Unsupervised:} Does not require labels. Methods include Isolation Forest, One-Class SVM, DBSCAN, and autoencoders. These methods learn the ``normal'' distribution and flag deviations.
\end{itemize}

\subsection{Isolation Forest}

Isolation Forest (IF), introduced by Liu et al. (2008), is particularly well-suited to telecom anomaly detection because:

\begin{enumerate}
  \item It operates on numeric tabular data at high speed and low memory.
  \item It does not require labels, which matches the real-world telecom scenario.
  \item It naturally handles high-dimensional feature spaces.
  \item Its ``contamination'' parameter allows domain-informed tuning of anomaly prevalence.
  \item It produces both hard decisions ($\pm 1$) and continuous anomaly scores.
\end{enumerate}

IF constructs an ensemble of random trees. Each tree recursively selects a random feature and a random split value. Anomalies --- being ``few and different'' --- require fewer splits to be isolated, resulting in shorter path lengths. The anomaly score is derived from the average path length across all trees.

\subsection{Gradient Boosting for Risk Prediction}

For supervised regression tasks such as SLA risk prediction, Gradient Boosting Regressors (GBR) are the method of choice for tabular telecom data. GBR builds an ensemble of weak learners (shallow decision trees) sequentially, with each tree correcting the errors of the previous ensemble. Key advantages for our use case:

\begin{itemize}
  \item Handles non-linear KPI interactions (e.g., high latency \emph{and} packet loss amplifies risk more than either alone).
  \item Provides feature importances for explainability --- critical during the PFE defence.
  \item Strong performance on small-to-medium tabular datasets (our data scale).
  \item Well-understood, academically defensible, and supported by scikit-learn.
\end{itemize}

\section{OSS--BSS Convergence}

\subsection{The Convergence Problem}

OSS and BSS data have historically been managed by separate teams, stored in separate systems, and analysed with separate tools. The TM Forum's ODA (Open Digital Architecture) initiative explicitly identifies OSS--BSS convergence as a strategic priority. Huawei's O+B convergence vision aligns with this broader industry movement.

\subsection{Statistical Correlation Methods}

We employ two complementary correlation methods:

\begin{itemize}
  \item \textbf{Pearson correlation:} Measures the \emph{linear} relationship between two variables. A Pearson coefficient of $r = -0.8$ between latency and revenue would indicate that as latency increases, revenue decreases linearly.
  \item \textbf{Spearman correlation:} Measures the \emph{monotonic} (rank-based) relationship. It captures non-linear but consistent trends --- for example, revenue tends to decrease when packet loss increases, even if the relationship is not strictly linear.
\end{itemize}

By computing both Pearson and Spearman on multiple OSS--BSS metric pairs, we produce a comprehensive picture of how network degradation manifests in business outcomes.

\section{Cloud-Native Architecture in Telecom}

\subsection{The Shift to Cloud-Native}

The telecom industry is undergoing a fundamental architectural shift from monolithic network functions to cloud-native, containerised microservices. This shift is driven by:

\begin{itemize}
  \item \textbf{ETSI NFV:} Network Functions Virtualisation decouples software from hardware.
  \item \textbf{5G SBA:} 5G's Service-Based Architecture defines network functions as loosely coupled services.
  \item \textbf{Kubernetes / Docker:} Container orchestration enables horizontal scaling, rolling updates, and infrastructure abstraction.
\end{itemize}

\subsection{Huawei Cloud Stack (HCS)}

Huawei Cloud Stack is Huawei's hybrid cloud platform, deployed on-premises at operator sites. It provides:

\begin{itemize}
  \item \textbf{OBS (Object Storage Service):} S3-compatible object storage --- equivalent to our MinIO.
  \item \textbf{RDS (Relational Database Service):} Managed PostgreSQL/MySQL --- equivalent to our PostgreSQL container.
  \item \textbf{ECS/CCE:} Elastic Cloud Server / Cloud Container Engine --- container runtime equivalent to our Docker Compose.
  \item \textbf{VPC:} Network isolation and security boundaries.
\end{itemize}

Our architecture is designed from inception for HCS portability: the same Docker containers, with environment variable changes, can run on HCS with minimal modification.

\section{Gap Analysis and Positioning}

% ── Figure: Gap Analysis ──────────────────────────────────────────────────────
\begin{figure}[H]
\centering
\begin{tikzpicture}[
  gapnode/.style={
    rectangle, rounded corners=4pt, minimum width=4.5cm, minimum height=1.5cm,
    font=\small, align=center, line width=0.7pt,
  },
]
  % Existing solutions
  \node[gapnode, draw=ossblue!60, fill=ossblue!10] (oss) at (0,0) {
    \textbf{OSS Analytics}\\
    Network anomaly detection\\
    KPI monitoring
  };
  \node[gapnode, draw=bssgreen!60, fill=bssgreen!10] (bss) at (6,0) {
    \textbf{BSS Analytics}\\
    Revenue analytics\\
    Churn prediction
  };
  
  % Gap
  \node[gapnode, draw=huaweired!60, fill=huaweired!8,
        minimum width=3cm, minimum height=1.5cm,
        dashed, line width=1.2pt] (gap) at (3,0) {
    \textbf{GAP}\\
    \scriptsize No correlation\\
    \scriptsize No unified agent
  };
  
  % Our solution
  \node[gapnode, draw=aiviolet!70, fill=aiviolet!12,
        minimum width=10.5cm, minimum height=1.5cm,
        line width=1.2pt] (ours) at (3,-2.5) {
    \textbf{Our AI Operations Agent}\\
    O+B Convergence $|$ Real data $|$ Cloud-native $|$ CEM $\to$ CVM
  };
  
  % Arrows
  \draw[myarrow=aiviolet!70, line width=1pt] (gap) -- (ours)
    node[arrlabel, midway, right, xshift=3pt] {We fill this};
    
\end{tikzpicture}
\caption{Gap analysis: our AI Operations Agent bridges the OSS--BSS analytics gap.}
\label{fig:gap-analysis}
\end{figure}

The state of the art leaves a clear architectural gap: no existing PFE-scale solution combines real operator data, O+B convergence, cloud-native deployment, and an agentic AI architecture that bridges CEM and CVM. Our project fills this gap.

\section{Chapter Summary}

This chapter surveyed AI in telecom operations, anomaly detection methods (Isolation Forest, Gradient Boosting), OSS--BSS convergence challenges, and cloud-native telecom architectures. We identified the gap between existing CEM and CVM platforms --- the absence of a unified AI intelligence layer --- and positioned our project as the solution. The next chapter details the system architecture.
